\documentclass[12pt, a4paper]{article}

% Поддержка русского языка и кодировок
\usepackage[utf8]{inputenc}
\usepackage[T2A]{fontenc}
\usepackage[russian]{babel}

% Математические пакеты
\usepackage{amsmath, amssymb, amsthm, amsfonts}

% Настройка полей страницы
\usepackage{geometry}
\geometry{left=2cm, right=2cm, top=2cm, bottom=2cm}

% Кликабельное оглавление
\usepackage{hyperref}
\hypersetup{
	colorlinks=true,
	linkcolor=blue,
	filecolor=magenta,      
	urlcolor=cyan,
}

% Настройка окружений для теорем, определений и доказательств
\theoremstyle{definition}
\newtheorem{definition}{Определение}[section]
\newtheorem{example}{Пример}[section]

\theoremstyle{plain}
\newtheorem{theorem}{Теорема}[section]
\newtheorem{lemma}{Лемма}[section]
\newtheorem{property}{Свойство}[section]

% Переопределение символа конца доказательства на закрашенный квадрат
\renewcommand{\qedsymbol}{$\blacksquare$}

\title{\textbf{Ответы на вопросы к экзамену по алгебре}}
\author{Весна 2026}
\date{}

\begin{document}
	
	\maketitle
	\tableofcontents
	\newpage
	
	\section{Определение кольца. Делители нуля и регулярные элементы. Доказать, обратимый элемент является регулярным. Обратимые элементы кольца. Доказать, что в кольце с единицей обратимые элементы образуют группу. Область целостности. Степени и кратные в кольце.}
	
	\begin{definition}
		Алгебраическая система $(R, +, \cdot)$ называется \textbf{кольцом}, если:
		\begin{enumerate}
			\item $(R, +)$ — абелева группа.
			\item Умножение ассоциативно: $(ab)c = a(bc)$.
			\item Выполнены законы дистрибутивности: $a(b+c) = ab+ac$ и $(a+b)c = ac+bc$.
		\end{enumerate}
	\end{definition}
	
	\begin{definition}
		Элемент $a \in R$ называется \textbf{делителем нуля}, если существует такой ненулевой элемент $b \neq 0$, что $ab=0$ или $ba=0$.
	\end{definition}
	
	\begin{definition}
		Элемент $a \in R$ называется \textbf{регулярным} (левым/правым), если на него можно сокращать равенства: $ab=ac \implies b=c$.
	\end{definition}
	
	\begin{theorem}
		Всякий обратимый элемент кольца является регулярным.
	\end{theorem}
	\begin{proof}
		Пусть $a$ — обратимый элемент, то есть существует $a^{-1}$. Если $ab = ac$, то, умножив обе части равенства слева на $a^{-1}$, получим $a^{-1}ab = a^{-1}ac \implies e \cdot b = e \cdot c \implies b = c$. Следовательно, $a$ — регулярный.
	\end{proof}
	
	\begin{theorem}
		В кольце с единицей обратимые элементы образуют группу по умножению.
	\end{theorem}
	\begin{proof}
		Пусть $U(R)$ — множество обратимых элементов.
		\begin{enumerate}
			\item Замкнутость: если $a, b \in U(R)$, то $(ab)^{-1} = b^{-1}a^{-1} \in R$, значит $ab \in U(R)$.
			\item Наличие нейтрального элемента: $1 \in U(R)$, так как $1 \cdot 1 = 1$.
			\item Ассоциативность: наследуется из аксиом кольца.
			\item Наличие обратного: для каждого $a \in U(R)$ существует $a^{-1}$, который также обратим, так как $(a^{-1})^{-1} = a$.
		\end{enumerate}
		Таким образом, $U(R)$ — группа.
	\end{proof}
	
	\begin{definition}
		\textbf{Областью целостности} называется коммутативное кольцо с единицей без делителей нуля.
	\end{definition}
	
	\textbf{Степени и кратные в кольце:} Для любого элемента $a \in R$ натуральная степень определяется как $a^n = a \cdot a \cdot \dots \cdot a$ ($n$ раз). Кратное элемента определяется через сложение: $na = a + a + \dots + a$ ($n$ раз). Для них справедливы стандартные свойства: $a^n a^m = a^{n+m}$ и $na + ma = (n+m)a$.
	
	\newpage
	\section{Определение поля. Мультипликативная группа поля. Регулярные элементы поля. Доказать, что поле является областью целостности. Определение подкольца, унитарного подкольца и подполя. Докажите характеристическое свойство подполя. Приведите примеры подколец, в которых единицы совпадают или отличаются от единиц в надкольце.}
	
	\begin{definition}
		\textbf{Полем} называется коммутативное кольцо с единицей $1 \neq 0$, в котором каждый ненулевой элемент обратим.
	\end{definition}
	
	\textbf{Мультипликативная группа поля:} Множество всех ненулевых элементов поля $K$, обозначаемое $K^* = K \setminus \{0\}$, образует абелеву группу по умножению.
	
	\textbf{Регулярные элементы поля:} Поскольку в поле каждый ненулевой элемент обратим, а обратимые элементы регулярны, то все ненулевые элементы поля являются регулярными.
	
	\begin{theorem}
		Всякое поле является областью целостности.
	\end{theorem}
	\begin{proof}
		Пусть $K$ — поле. По определению оно коммутативно и содержит $1 \neq 0$. Пусть $a, b \in K$ и $ab = 0$. Если $a \neq 0$, то существует $a^{-1}$. Умножим равенство на $a^{-1}$ слева: $a^{-1}ab = a^{-1} \cdot 0 \implies 1 \cdot b = 0 \implies b = 0$. Делителей нуля нет, значит $K$ — область целостности.
	\end{proof}
	
	\begin{definition}
		Непустое подмножество $S$ кольца $R$ называется \textbf{подкольцом}, если оно само является кольцом относительно индуцированных операций. Эквивалентно: для любых $a,b \in S$ выполняется $a-b \in S$ и $a \cdot b \in S$.
	\end{definition}
	
	\begin{definition}
		Если исходное кольцо $R$ содержит единицу $1_R$, то подкольцо $S \subseteq R$ называется \textbf{унитарным подкольцом}, если оно содержит ту же самую единицу: $1_S = 1_R$.
	\end{definition}
	
	\begin{definition}
		Непустое подмножество $P$ поля $K$ называется \textbf{подполем}, если $P$ само является полем относительно операций, определенных в $K$.
	\end{definition}
	
	\begin{theorem}[Характеристическое свойство подполя]
		Подмножество $P$ поля $K$ является подполем тогда и только тогда, когда оно содержит не менее двух элементов и замкнуто относительно вычитания и деления на ненулевой элемент.
	\end{theorem}
	\begin{proof}
		Необходимость очевидна из аксиом поля. \\
		Достаточность: Пусть свойства выполнены. Так как $P$ содержит не менее двух элементов, в нем есть $a \neq 0$. Тогда $a-a = 0 \in P$. Далее, $0-a = -a \in P$. Замкнутость по сложению: $a - (-b) = a+b \in P$. Для умножения: $a \cdot (b^{-1})^{-1}$ выводится через замкнутость к делению $a/b \in P$ при $b \neq 0$.
	\end{proof}
	
	\begin{example}Примеры соотношения единиц:
		\begin{enumerate}
			\item \textbf{Единицы совпадают:} Кольцо $\mathbb{Z}$ является унитарным подкольцом в поле $\mathbb{Q}$. Их единицы совпадают: $1_{\mathbb{Z}} = 1_{\mathbb{Q}} = 1$.
			\item \textbf{В подкольце нет единицы:} Кольцо четных чисел $2\mathbb{Z}$ является подкольцом в $\mathbb{Z}$, но не имеет единичного элемента.
			\item \textbf{Единицы отличаются:} В кольце $R = \mathbb{Z} \times \mathbb{Z}$ единицей является пара $(1,1)$. Подмножество $S = \{(a,0) \mid a \in \mathbb{Z}\}$ — подкольцо, единицей которого является $(1,0)$. Очевидно, что $1_S \neq 1_R$.
		\end{enumerate}
	\end{example}
	
	\newpage
	\section{Гомоморфизм колец. Докажите теорему о 4-х свойствах гомоморфизма колец. Докажите характеристическое свойство изоморфизма колец. Докажите, что ядро гомоморфизма - идеал.}
	
	\begin{definition}
		Отображение $\phi: R \to S$ называется \textbf{гомоморфизмом колец}, если для любых $a, b \in R$ выполняются условия: $\phi(a+b) = \phi(a) + \phi(b)$ и $\phi(ab) = \phi(a)\phi(b)$.
	\end{definition}
	
	\begin{theorem}[Четыре свойства гомоморфизма]
		Пусть $\phi: R \to S$ — гомоморфизм колец. Тогда:
		\begin{enumerate}
			\item $\phi(0_R) = 0_S$.
			\item $\phi(-a) = -\phi(a)$.
			\item $\phi(a-b) = \phi(a)-\phi(b)$.
			\item Если $R'$ — подкольцо в $R$, то $\phi(R')$ — подкольцо в $S$.
		\end{enumerate}
	\end{theorem}
	\begin{proof}
		1) $\phi(0) = \phi(0+0) = \phi(0) + \phi(0)$. Сокращая на $\phi(0)$ в группе $(S, +)$, получаем $\phi(0) = 0$. \\
		2) $\phi(a) + \phi(-a) = \phi(a + (-a)) = \phi(0) = 0 \implies \phi(-a) = -\phi(a)$. \\
		3) Очевидно следует из пунктов 1 и 2: $\phi(a + (-b)) = \phi(a) + \phi(-b) = \phi(a) - \phi(b)$. \\
		4) Проверяется через замкнутость операций сложения и умножения на образе отображения $\phi(R')$.
	\end{proof}
	
	\textbf{Характеристическое свойство изоморфизма колец:}
	Изоморфизм — это биективный гомоморфизм. Характеристическое свойство заключается в том, что отображение $\phi: R \to S$ является изоморфизмом тогда и только тогда, когда существует обратное отображение $\phi^{-1}: S \to R$, которое также является гомоморфизмом колец. Доказывается проверкой сохранения операций при применении обратного отображения.
	
	\begin{theorem}
		Ядро гомоморфизма $\ker \phi = \{a \in R \mid \phi(a) = 0\}$ является двусторонним идеалом кольца $R$.
	\end{theorem}
	\begin{proof}
		Пусть $x, y \in \ker \phi$ и $r \in R$.
		\begin{enumerate}
			\item $\phi(x-y) = \phi(x) - \phi(y) = 0 - 0 = 0 \implies x-y \in \ker \phi$.
			\item $\phi(rx) = \phi(r)\phi(x) = \phi(r) \cdot 0 = 0 \implies rx \in \ker \phi$.
			\item $\phi(xr) = \phi(x)\phi(r) = 0 \cdot \phi(r) = 0 \implies xr \in \ker \phi$.
		\end{enumerate}
		Поскольку $\ker \phi$ является подгруппой по сложению и выдерживает умножение на элементы из $R$ как слева, так и справа, то это двусторонний идеал.
	\end{proof}
	
	\newpage
	\setcounter{section}{3}
	\section{Определение делимости в кольцах. Обратимые элементы. Ассоциированные элементы в области целостности. Несобственные делители. Неразложимые элементы. Область с однозначным разложением на множители. Примеры. Теорема о двух эквивалентных определениях ООР.}
	
	\begin{definition}
		Пусть $R$ — область целостности. Говорят, что элемент $a \in R$ \textbf{делит} элемент $b \in R$ (обозначение $a \mid b$), если существует $c \in R$ такое, что $b=ac$.
	\end{definition}
	
	\begin{definition}
		Элемент $u \in R$ называется \textbf{обратимым}, если существует элемент $v \in R$ такой, что $uv=vu=1$. Элемент $v$ называется обратным к $u$ и обозначается $u^{-1}$.
	\end{definition}
	
	\begin{definition}
		В области целостности элементы $a$ и $b$ называются \textbf{ассоциированными} ($a \sim b$), если $a = bu$, где $u$ — обратимый элемент. Отношение ассоциированности является отношением эквивалентности (рефлексивно, симметрично, транзитивно).
	\end{definition}
	
	\begin{definition}
		Делитель элемента $a$ называется \textbf{несобственным}, если он ассоциирован с $a$ или является обратимым элементом.
	\end{definition}
	
	\begin{definition}
		Ненулевой необратимый элемент $p$ называется \textbf{неразложимым}, если из равенства $p=ab$ следует, что либо $a$, либо $b$ обратим. Элемент $p$ неразложим тогда и только тогда, когда у него нет собственных делителей.
	\end{definition}
	
	\begin{definition}
		Область целостности называется \textbf{областью с однозначным разложением} (ООР, факториальным кольцом), если любой ненулевой необратимый элемент раскладывается в произведение неразложимых, и это разложение единственно с точностью до порядка множителей и ассоциированности.
	\end{definition}
	
	\begin{example}Примеры ООР:
		\begin{enumerate}
			\item Кольцо целых чисел $\mathbb{Z}$ является ООР. Например, $30 = 2 \cdot 3 \cdot 5$. Любое другое разложение отличается только порядком и умножением на $\pm 1$ .
			\item Кольцо многочленов $K[x]$ над полем $K$ является ООР.
		\end{enumerate}
	\end{example}
	
	\begin{theorem}[Два эквивалентных определения ООР]
		Для области целостности $R$ следующие определения ООР эквивалентны:
		1. Каждый ненулевой необратимый элемент раскладывается в произведение неразложимых, и разложение единственно с точностью до порядка и ассоциированности.
		2. Каждый ненулевой необратимый элемент раскладывается в произведение неразложимых, и всякий неразложимый элемент является простым.
	\end{theorem}
	\begin{proof}
		\textbf{(1) $\Rightarrow$ (2):} Пусть $p$ — неразложимый элемент и $p \mid ab$. Тогда $ab = pc$ для некоторого $c$. Разложим $a, b, c$ и $p$ на неразложимые множители. В силу единственности разложения ООР, $p$ должен быть ассоциирован с одним из множителей произведения $ab$. Следовательно, $p$ ассоциирован либо с множителем из $a$, либо из $b$, откуда $p \mid a$ или $p \mid b$. Значит, $p$ прост.
		
		\textbf{(2) $\Rightarrow$ (1):} Существование разложения предполагается. Докажем единственность. Пусть $p_1 \dots p_n = q_1 \dots q_m$, где все множители неразложимы. Так как $p_1$ прост и делит правую часть, он должен делить какой-то $q_j$. Но $q_j$ неразложим, поэтому $p_1 \sim q_j$. Переставляя множители и сокращая на $p_1$ (вынося обратимый элемент $u$), можно рекурсивно поглотить все множители. В итоге получаем $n=m$ и $p_i \sim q_i$ (с точностью до перестановки). Разложение единственно.
	\end{proof}
	
	\newpage
	\section{Свойства взаимно простых элементов. Определение евклидовой области. Примеры. Следствия из определения евклидовой области (4 следствия).}
	
	\begin{definition}
		Элементы $a,b \in R$ называются \textbf{взаимно простыми}, если их наибольший общий делитель обратим: $(a,b)=1$.
	\end{definition}
	
	\textbf{Свойства взаимно простых элементов:}
	\begin{enumerate}
		\item \textit{Соотношение Безу:} $(a,b)=1 \iff$ существуют $x, y \in R$ такие, что $xa+yb=1$ .
		\item Если элементы $a$ и $b$ взаимно просты и $a \mid bc$, то $a \mid c$.
		\item Если элемент $a$ взаимно прост с каждым из элементов $b_1, \dots, b_n$, то он взаимно прост с их произведением $b_1 \dots b_n$ .
	\end{enumerate}
	
	\begin{definition}
		Область целостности $R$ называется \textbf{евклидовой}, если существует функция $N: R \setminus \{0\} \to \mathbb{Z}_{\ge 0}$ (евклидова норма), такая что :
		1. Для любых ненулевых $a,b \in R$: $N(a) \le N(ab)$.
		2. Для любых $a, b \in R$ ($b \neq 0$) существуют элементы $q, r \in R$ такие, что $a = bq + r$, где либо $r=0$, либо $N(r) < N(b)$ (деление с остатком) .
	\end{definition}
	
	\begin{example}Примеры евклидовых областей:
		\begin{enumerate}
			\item Кольцо целых чисел $\mathbb{Z}$ с нормой $N(a) = |a|$.
			\item Кольцо многочленов $K[x]$ над полем $K$ с нормой $N(f) = \deg f$.
			\item Кольцо гауссовых целых чисел $\mathbb{Z}[i]$ с нормой $N(a+bi) = a^2+b^2$.
		\end{enumerate}
	\end{example}
	
	\begin{theorem}[Следствия из определения евклидовой области]
		Для любой евклидовой области справедливы следующие 4 утверждения:
		\begin{enumerate}
			\item \textbf{Существует алгоритм Евклида} для нахождения НОД. Последовательность делений с остатком обрывается, так как нормы остатков строго убывают: $N(b) > N(r_1) > N(r_2) > \dots$. Последний ненулевой остаток — это НОД.
			\item \textbf{Любой идеал является главным}. Для ненулевого идеала $I$ выбирается элемент $d \in I$ с минимальной нормой. Делением с остатком любого $a \in I$ на $d$ показывается, что остаток $r=0$, следовательно $a=dq$ и $I=(d)$.
			\item \textbf{Элемент имеет минимальную норму $\iff$ он обратим}.
			\item \textbf{Любые два элемента имеют НОД}, представимый в виде линейной комбинации $d = xa+yb$. Множество $I = \{xa+yb \mid x,y \in R\}$ является идеалом, а значит, порождается одним элементом $d$, который и является НОД.
		\end{enumerate}
	\end{theorem}
	
	\newpage
	\section{Доказать, что евклидова область является ООР.}
	
	\begin{theorem}
		Всякая евклидова область является областью с однозначным разложением (ООР).
	\end{theorem}
	\begin{proof}
		Полноценное доказательство состоит из двух шагов:
		\begin{enumerate}
			\item \textbf{Существование разложения:} Любой необратимый ненулевой элемент можно разложить на множители. Если элемент разложим, мы делим его на собственные множители. Поскольку при делении норма строго уменьшается ($N(a) < N(ab)$ для необратимых элементов), цепочка собственных делителей не может быть бесконечной. Следовательно, процесс оборвется на неразложимых элементах.
			\item \textbf{Единственность разложения:} Базируется на том факте, что каждый неразложимый элемент в евклидовом кольце является простым. А именно: если неразложимый элемент $p$ делит произведение $ab$, то он делит $a$ или $b$. Это доказывается через линейное представление НОД в евклидовых кольцах (соотношение Безу, следствие 4 из предыдущего билета). Опираясь на простоту неразложимых элементов, единственность разложения доказывается стандартной индукцией по числу множителей (аналогично второй части теоремы об эквивалентных определениях ООР).
		\end{enumerate}
	\end{proof}
	
	\newpage
	\setcounter{section}{6}
	\section{Полиномы от одной неизвестной. Проверить выполнение аксиом кольца, наличие единицы. Доказать вложение кольца коэффициентов в кольцо полиномов. Доказать, что если R- область целостности, то R[x] тоже. Как представляются одночлены финитными последовательностями в R[x]?}
	
	\begin{definition}
		Множество финитных последовательностей $(a_0, a_1, \dots, a_n, 0, \dots)$ элементов кольца $R$ с покомпонентным сложением и умножением по правилу Коши $c_k = \sum_{i+j=k} a_i b_j$ называется \textbf{кольцом полиномов} $R[x]$ .
	\end{definition}
	
	\textbf{Представление одночленов:} Одночлен $x^n$ представляется финитной последовательностью $(0, \dots, 0, 1, 0, \dots)$, где единица стоит на $n$-й позиции .
	
	\begin{theorem}
		Множество $R[x]$ является кольцом. Если $R$ имеет единицу, то $R[x]$ также имеет единицу.
	\end{theorem}
	\begin{proof}
		Аксиомы кольца проверяются напрямую . Сумма финитных последовательностей финитна. Ассоциативность, коммутативность сложения и дистрибутивность наследуются из $R$. Если в $R$ есть $1$, то единицей в $R[x]$ будет последовательность $1_{R[x]} = (1, 0, 0, \dots)$ .
	\end{proof}
	
	\begin{theorem}
		Кольцо $R$ вкладывается в кольцо $R[x]$.
	\end{theorem}
	\begin{proof}
		Отображение $\varphi(a) = (a, 0, 0, \dots)$ является инъективным гомоморфизмом, так как сохраняет сложение и умножение, и $\varphi(a) = \varphi(b) \implies a=b$ .
	\end{proof}
	
	\begin{theorem}
		Если $R$ — область целостности, то $R[x]$ — тоже.
	\end{theorem}
	\begin{proof}
		Пусть $f, g \neq 0$. Их старшие коэффициенты $a_n \neq 0$ и $b_m \neq 0$. Коэффициент при старшей степени $x^{n+m}$ в произведении $f \cdot g$ равен $a_n b_m$. Так как $R$ — область целостности, $a_n b_m \neq 0$, следовательно, $f \cdot g \neq 0$, то есть делителей нуля нет .
	\end{proof}
	
	\newpage
	
	\section{Доказать, что отображение подстановки в полином элемента из унитарного надкольца коэффициентов является гомоморфизмом кольца полиномов в это надкольцо. Определение алгебраического и трансцендентного элементов. Доказать, что все кольца полиномов над R изоморфны.}
	
	\begin{theorem}
		Отображение $\phi_x: R[t] \to S$, заданное как $\phi_x(f(t)) = f(x)$, является гомоморфизмом.
	\end{theorem}
	\begin{proof}
		Проверяем сохранение сложения: $\phi_x(f+g) = (f+g)(x) = f(x)+g(x) = \phi_x(f) + \phi_x(g)$ . Сохранение умножения вытекает из того, что $(fg)(x) = \sum (\sum a_i b_j) x^k = f(x)g(x)$ .
	\end{proof}
	
	\begin{definition}
		Элемент $x \in S$ \textbf{алгебраичен} над $R$, если существует ненулевой $f \in R[t]$ такой, что $f(x)=0$ ($\ker \phi_x \neq \{0\}$). Если $\ker \phi_x = \{0\}$, элемент \textbf{трансцендентен} .
	\end{definition}
	
	\begin{theorem}
		Все кольца полиномов над $R$ изоморфны.
	\end{theorem}
	\begin{proof}
		Отображение $\varphi: R[x] \to R[y]$, заданное как $\varphi(\sum a_i x^i) = \sum a_i y^i$, биективно и сохраняет сложение и умножение, так как правила операций с коэффициентами не зависят от обозначения переменной .
	\end{proof}
	
	\newpage
	
	\section{Доказать, что конечная область целостности является полем. Доказать, что если кольцо можно вложить в поле, то оно является областью целостности, и обратно.}
	
	\begin{theorem}
		Всякая конечная область целостности является полем.
	\end{theorem}
	\begin{proof}
		Пусть $R = \{x_1, \dots, x_n\}$. Для любого $a \neq 0$ отображение $f(x) = ax$ инъективно, так как $ax_i = ax_j \implies a(x_i - x_j) = 0 \implies x_i = x_j$ (нет делителей нуля). Поскольку множество конечно, инъективность влечет сюръективность. Значит, существует $b \in R$, для которого $ab=1$, то есть $a$ обратим .
	\end{proof}
	
	\begin{theorem}
		Если кольцо вкладывается в поле, то оно является областью целостности.
	\end{theorem}
	\begin{proof}
		Пусть есть инъективный гомоморфизм $\varphi: R \to K$ (где $K$ — поле). Если $ab=0$, то $\varphi(a)\varphi(b)=0$. В поле нет делителей нуля, значит $\varphi(a)=0$ или $\varphi(b)=0$. Из инъективности следует $a=0$ или $b=0$ .
	\end{proof}
	
	\newpage
	
	\section{Построение поля частных для области целостности. Примеры полей частных: рациональные, гауссовы.}
	
	\textbf{Построение:} Для области целостности $R$ рассматривается множество пар $(a, b)$, где $a \in R, b \in R \setminus \{0\}$. Вводится отношение эквивалентности: $(a,b) \sim (c,d) \iff ad=bc$ . Класс эквивалентности обозначается как дробь $\frac{a}{b}$. Операции задаются стандартно: $\frac{a}{b} + \frac{c}{d} = \frac{ad+bc}{bd}$ и $\frac{a}{b} \cdot \frac{c}{d} = \frac{ac}{bd}$ . Это множество образует поле $Frac(R)$.
	
	\begin{example}
		Полем частных для $\mathbb{Z}$ является поле рациональных чисел $\mathbb{Q}$ . Полем частных для гауссовых целых чисел $\mathbb{Z}[i]$ является поле гауссовых рациональных чисел $\mathbb{Q}(i)$ .
	\end{example}
	
	\newpage
	
	\section{Поле рациональных дробей F(x). Доказать, что оно бесконечно. Что можно сказать о его характеристике?}
	
	\begin{theorem}
		Поле $F(x)$ бесконечно и его характеристика совпадает с характеристикой $F$.
	\end{theorem}
	\begin{proof}
		Оно содержит бесконечную последовательность элементов $x, x^2, x^3, \dots$, которые попарно различны (многочлены разной степени не равны) . Характеристика $F(x)$ совпадает с $char F$, так как единицы полей совпадают, и если $p \cdot 1_F = 0$, то и $p \cdot 1_{F(x)} = 0$.
	\end{proof}
	
	\newpage
	
	\section{Делители единицы, обратимые элементы, собственные делители элемента и ассоциированные элементы области целостности. НОД и НОК в области целостности.}
	
	В кольце многочленов $K[x]$ обратимыми элементами (делителями единицы) являются только ненулевые константы $K^*$. Ассоциированные элементы отличаются на обратимый множитель (константу) .
	
	\textbf{НОД} многочленов $f$ и $g$ — это такой их общий делитель $d$, который делится на любой другой их общий делитель . 
	\textbf{НОК} — это их общее кратное $L$, которое делит любое другое их общее кратное .
	Справедливо: $gcd(f,g) \cdot lcm(f,g) = f_0(x) \cdot g_0(x)$ (для унитарных представителей) .
	
	\newpage
	
	\section{Евклидовы кольца: определение, три примера с обоснованием.}
	
	\begin{definition}
		Область целостности $R$ называется евклидовой, если задана норма $N: R \setminus \{0\} \to \mathbb{N}_0$, и для любых $a, b \neq 0$ существует деление с остатком: $a = qb + r$, где $r=0$ или $N(r) < N(b)$ .
	\end{definition}
	
	\begin{example}
		1. Кольцо $\mathbb{Z}$ с нормой $N(n) = |n|$ (обычное деление с остатком). \\
		2. $F[x]$ с нормой $N(f) = \deg f$ (деление многочленов "углом"). \\
		3. Гауссовы целые $\mathbb{Z}[i]$ с нормой $N(a+bi) = a^2+b^2$ (деление осуществляется через приближение дроби в $\mathbb{Q}(i)$) .
	\end{example}
	
	\newpage
	
	\section{Взаимно простые элементы, алгоритм Евклида и представление НОД в евклидовом кольце. Свойства взаимно простых элементов и произведений.}
	
	Алгоритм Евклида заключается в цепочке делений: $a = q_0 b + r_1$, $b = q_1 r_1 + r_2 \dots$. Так как нормы остатков строго убывают, процесс конечен. Последний ненулевой остаток — это НОД . Обратным ходом алгоритма НОД можно выразить как линейную комбинацию: $d = ua+vb$ .
	
	\begin{property}
		1. Если $a \perp b$ и $a \mid bc$, то $a \mid c$. \\
		2. Если $a \perp b$ и $a \perp c$, то $a \perp bc$.
	\end{property}
	
	\newpage
	
	\section{Неприводимые элементы области целостности. Примеры. Свойства: делимость произведения нескольких элементов на неприводимый. Теорема о факториальности евклидовых колец.}
	
	Ненулевой необратимый $p \in R$ неприводим, если из $p=ab$ следует, что $a$ или $b$ обратим. Примеры: простые числа в $\mathbb{Z}$, неприводимые многочлены в $F[x]$, $1+i$ в $\mathbb{Z}[i]$.
	
	\begin{theorem}
		В евклидовой области если неприводимый $p \mid a_1 \dots a_n$, то $p \mid a_i$ для некоторого $i$.
	\end{theorem}
	\begin{theorem}
		Каждая евклидова область факториальна (ООР) .
	\end{theorem}
	
	\newpage
	
	\section{Определение идеала. Примеры. Идеал с обратимым элементом. Радикал коммутативного кольца. Доказать, что он является идеалом.}
	
	\begin{definition}
		Подмножество $I \subseteq R$ — \textbf{идеал}, если оно подгруппа по сложению и поглощает умножение на элементы из $R$: $\forall r \in R, a \in I \implies ra, ar \in I$ .
	\end{definition}
	
	Примеры: $n\mathbb{Z}$ в $\mathbb{Z}$, главный идеал $(f)$ в $F[x]$.
	
	\begin{theorem}
		Если $I$ содержит обратимый $u$, то $I=R$.
	\end{theorem}
	\begin{proof}
		Так как $u \in I$ и идеал поглощает умножение, $u \cdot u^{-1} = 1 \in I$. Значит $\forall r \in R$, $r \cdot 1 = r \in I \implies I=R$ .
	\end{proof}
	
	\textbf{Радикал} $Rad(R)$ — это пересечение всех максимальных идеалов. Любое пересечение идеалов является идеалом, следовательно, $Rad(R)$ — идеал .
	
	\newpage
	
	\section{Главные идеалы коммутативного кольца. Примеры с обоснованием. Верно ли, что подкольцо является идеалом? Примеры. При каких условиях это верно?}
	
	\textbf{Главный идеал} $(a) = Ra = \{ra \mid r \in R\}$ порождается одним элементом . Пример: $(n) = n\mathbb{Z}$ в $\mathbb{Z}$.
	
	\begin{theorem}
		Подкольцо не обязано быть идеалом.
	\end{theorem}
	\begin{proof}
		Контрпример: $\mathbb{Z}$ — подкольцо в $\mathbb{Q}$, но не идеал. $2 \in \mathbb{Z}$, $\frac{1}{2} \in \mathbb{Q}$, однако $\frac{1}{2} \cdot 3 = \frac{3}{2} \notin \mathbb{Z}$ . Подкольцо $S$ будет идеалом только если оно поглощает умножение из $R$.
	\end{proof}
	
	\newpage
	
	\section{Доказать, что поле является простым кольцом. И наоборот.}
	
	Простое кольцо не имеет нетривиальных двусторонних идеалов.
	
	\begin{theorem}
		Коммутативное кольцо с единицей является полем $\iff$ оно простое.
	\end{theorem}
	\begin{proof}
		($\Rightarrow$) Пусть $F$ — поле, $I \neq 0$ — идеал. В нем есть $a \neq 0$, который обратим. Значит $1 \in I \implies I=F$. \\
		($\Leftarrow$) Пусть $R$ — простое кольцо, $a \neq 0$. Главный идеал $Ra \neq \{0\}$, значит $Ra=R$. Следовательно, существует $b \in R$, для которого $ba=1$, то есть $a$ обратим.
	\end{proof}
	
	\newpage
	
	\section{Доказать, что кольцо квадратных матриц над полем является простым.}
	
	\begin{theorem}
		Кольцо $M_n(F)$ простое.
	\end{theorem}
	\begin{proof}
		Пусть $I \neq \{0\}$ — идеал, $A \in I$ и $A_{ij} \neq 0$. Умножая матрицу $A$ на матричные единицы, получаем $E_{ki} A E_{jl} = A_{ij} E_{kl} \in I$. Так как $A_{ij}$ обратим в $F$, все базисные матрицы $E_{kl} \in I$, откуда $I = M_n(F)$ .
	\end{proof}
	
	\newpage
	
	\section{Докажите, что произведение идеалов является идеалом. Примеры для целых чисел и многочленов.}
	
	\begin{definition}
		Произведение $AB = \{ \sum_{i=1}^n a_i b_i \mid a_i \in A, b_i \in B \}$.
	\end{definition}
	
	\begin{theorem}
		$AB$ — идеал.
	\end{theorem}
	\begin{proof}
		Замкнуто по сложению по определению (суммы сумм). Для $r \in R$, $r \sum a_i b_i = \sum (ra_i)b_i \in AB$, так как $A$ — идеал и $ra_i \in A$ .
	\end{proof}
	
	Примеры: В $\mathbb{Z}$: $(6)(10) = (60)$. В $F[x]$: $(f)(g) = (fg)$.
	
	\newpage
	\setcounter{section}{20}
	\section{Сумма идеалов. Её свойства, примеры в целых числах и многочленах. Доказать дистрибутивность умножения идеалов относительно сложения.}
	
	\begin{definition}
		Суммой идеалов $A$ и $B$ называется множество $A+B = \{a+b \mid a \in A, b \in B\}$ .
	\end{definition}
	
	\textbf{Свойства:} $A+B$ является наименьшим идеалом, содержащим $A$ и $B$ . Также справедливы включения $A \subseteq A+B$ и $B \subseteq A+B$. Операция сложения идеалов коммутативна и ассоциативна, а также $A+A=A$ .
	
	\textbf{Примеры:} В кольце $\mathbb{Z}$: $(m)+(n) = (\gcd(m,n))$, например $(6)+(10)=(2)$. В $F[x]$: $(f)+(g) = (\gcd(f,g))$.
	
	\begin{theorem}[Дистрибутивность]
		Для идеалов $A, B, C$ коммутативного кольца верно: $A(B+C) = AB+AC$ .
	\end{theorem}
	\begin{proof}
		($\supseteq$) Типичный порождающий элемент суммы $AB+AC$ имеет вид $ab+ac = a(b+c)$, что очевидно принадлежит $A(B+C)$. \\
		($\subseteq$) Произвольный элемент из $AB+AC$ имеет вид $\sum a_i b_i + \sum a'_j c'_j$. Каждое слагаемое в этой сумме — это произведение элемента из $A$ ($a_i$ или $a'_j$) на элемент из $B$ или $C$ (значит, из $B+C$). Следовательно, вся сумма по определению лежит в идеале $A(B+C)$ .
	\end{proof}
	
	\newpage
	
	\section{Доказать, что пересечение идеалов - идеал. Образует ли множество идеалов кольца кольцо? Примеры в целых числах и многочленах. Докажите, что произведение идеалов содержится в их пересечении.}
	
	\begin{theorem}
		Если $A, B$ — идеалы кольца $R$, то $A \cap B$ также является идеалом.
	\end{theorem}
	\begin{proof}
		Так как $0 \in A$ и $0 \in B$, то $0 \in A \cap B$. Если $x,y \in A \cap B$, то $x,y \in A$ и $x,y \in B$. Поскольку $A$ и $B$ — идеалы, их разность $x-y \in A$ и $x-y \in B$, значит $x-y \in A \cap B$ . Для любого $r \in R$, элементы $rx, xr \in A$ и $rx, xr \in B$, поэтому $rx, xr \in A \cap B$ . Следовательно, $A \cap B$ — идеал.
	\end{proof}
	
	\textbf{Образует ли множество идеалов кольцо?} Обычно нет. Относительно операций суммы и пересечения структура идеалов не образует кольцо, так как нет обратных элементов по сложению (идеал нельзя "вычесть"). Например, в $\mathbb{Z}$: $(2)+(3)=\mathbb{Z}$, $(2) \cap (3) = (6)$ .
	
	\begin{theorem}
		Для любых идеалов $A, B$ коммутативного кольца верно: $AB \subseteq A \cap B$ .
	\end{theorem}
	\begin{proof}
		Любой элемент из $AB$ имеет вид $\sum a_i b_i$, где $a_i \in A, b_i \in B$ . Так как $A$ — идеал, произведение $a_i b_i \in A$. Аналогично, так как $B$ — идеал, $a_i b_i \in B$ . Следовательно, вся конечная сумма лежит в $A \cap B$.
	\end{proof}
	
	\newpage
	
	\section{Взаимно простые идеалы в коммутативном кольце с единицей. Примеры в целых числах и многочленах. Доказать, что идеалы А, В взаимно просты тогда и только тогда, когда для некоторых их элементов a, b верно равенство: a+b=e.}
	
	\begin{definition}
		Идеалы $A, B$ коммутативного кольца с единицей называются \textbf{взаимно простыми}, если $A+B=R$ .
	\end{definition}
	
	\textbf{Примеры:} В кольце $\mathbb{Z}$ идеалы $(m)$ и $(n)$ взаимно просты тогда и только тогда, когда $\gcd(m,n)=1$. В кольце многочленов аналогично: $(f)$ и $(g)$ взаимно просты, если многочлены взаимно просты ($\gcd=1$).
	
	\begin{theorem}
		Идеалы $A, B$ взаимно просты $\iff$ существуют $a \in A, b \in B$ такие, что $a+b=e$ (где $e$ — единица кольца) .
	\end{theorem}
	\begin{proof}
		($\Rightarrow$) Если $A+B=R$, то единица $e \in A+B$, следовательно, $e = a+b$ для некоторых $a \in A, b \in B$ . \\
		($\Leftarrow$) Если существует разложение $e=a+b$, то для любого $r \in R$ имеем $r = re = r(a+b) = ra+rb$ . Так как $ra \in A$ и $rb \in B$, то $r \in A+B$. Значит, $A+B=R$.
	\end{proof}
	
	\newpage
	
	\section{Доказать, что если идеал взаимно прост с несколькими идеалами, то он взаимно прост и с их произведением. Идеал, порождённый данным конечным множеством элементов. В каком виде могут быть представлены любые его элементы?}
	
	\begin{theorem}
		Если идеал $A$ взаимно прост с идеалами $B_1, \dots, B_n$, то он взаимно прост с их произведением $B_1 \dots B_n$ .
	\end{theorem}
	\begin{proof}
		Достаточно рассмотреть случай двух идеалов. Пусть $A+B=R$ и $A+C=R$ . Тогда существуют $a_1, a_2 \in A$, $b \in B$, $c \in C$ такие, что $a_1+b=e$ и $a_2+c=e$ . Перемножая их, получаем: $e = (a_1+b)(a_2+c) = a_1 a_2 + a_1 c + b a_2 + bc$. Первые три слагаемых лежат в идеале $A$, а $bc \in BC$. Следовательно, $e \in A+BC$, то есть $A+BC=R$ .
	\end{proof}
	
	\textbf{Идеал, порождённый конечным множеством элементов} $a_1, \dots, a_n$, обозначается $(a_1, \dots, a_n)$ .
	Любые его элементы могут быть представлены в виде линейной комбинации с коэффициентами из кольца: $r_1 a_1 + \dots + r_n a_n$, где $r_i \in R$ .
	
	\newpage
	
	\section{Определение простого и максимального идеалов. Доказать, что нулевой идеал коммутативного кольца с единицей прост тогда и тогда, когда это кольцо область целостности.}
	
	\begin{definition}
		Идеал $P \neq R$ коммутативного кольца называется \textbf{простым}, если из $ab \in P$ следует $a \in P$ или $b \in P$ .
	\end{definition}
	
	\begin{definition}
		Идеал $M \neq R$ называется \textbf{максимальным}, если между $M$ и $R$ нет промежуточных идеалов.
	\end{definition}
	
	\begin{theorem}
		Нулевой идеал $(0)$ прост $\iff$ кольцо является областью целостности.
	\end{theorem}
	\begin{proof}
		Идеал $(0)$ прост тогда и только тогда, когда из $ab \in (0)$ (то есть $ab=0$) следует $a \in (0)$ или $b \in (0)$ (то есть $a=0$ или $b=0$) . Это в точности является определением области целостности (отсутствие делителей нуля).
	\end{proof}
	
	\newpage
	
	\section{Докажите, что каждый максимальный идеал прост.}
	
	\begin{theorem}
		Каждый максимальный идеал коммутативного кольца с единицей является простым.
	\end{theorem}
	\begin{proof}
		Пусть $M$ — максимальный идеал, и пусть $ab \in M$, но $a \notin M$ . Рассмотрим идеал $M+(a)$ . Так как $M$ максимален и $a \notin M$, то $M+(a) = R$ . Следовательно, существуют $m \in M$ и $r \in R$ такие, что $m+ra=e$ . Умножим обе части на $b$: $mb+rab=b$ . Так как $mb \in M$ и $rab = r(ab) \in M$ (ведь $ab \in M$), сумма $mb+rab$ принадлежит $M$. Значит, $b \in M$. Следовательно, $M$ прост.
	\end{proof}
	
	\newpage
	
	\section{Докажите, что любой необратимый элемент кольца содержится в максимальном идеале. Какие идеалы являются простыми в кольце целых чисел?}
	
	\begin{theorem}
		Любой необратимый элемент кольца с единицей содержится в некотором максимальном идеале.
	\end{theorem}
	\begin{proof}
		Пусть $a$ — необратимый элемент. Тогда порожденный им главный идеал $(a) \neq R$. Рассмотрим множество всех собственных идеалов, содержащих $(a)$. По лемме Цорна в этом множестве существует максимальный элемент $M$, который является максимальным идеалом всего кольца и содержит $a$.
	\end{proof}
	
	\textbf{Простые идеалы в $\mathbb{Z}$:} В кольце целых чисел простыми идеалами являются нулевой идеал $(0)$ и идеалы, порожденные простыми числами $(p)$ .
	
	\newpage
	
	\section{Докажите, что радикал коммутативного кольца совпадает с пересечением всех простых идеалов.}
	
	\begin{definition}
		\textbf{Радикалом} коммутативного кольца называется множество $Rad(R) = \{a \in R \mid a^n=0 \text{ для некоторого } n \ge 1\}$.
	\end{definition}
	
	\begin{theorem}
		$Rad(R) = \bigcap_{P \text{ - простой}} P$ .
	\end{theorem}
	\begin{proof}
		($\subseteq$) Если $a^n=0$, то $0 \in P$ для любого простого идеала $P$. Из определения простого идеала следует, что $a \in P$. Значит, $a \in \bigcap P_i$. \\
		($\supseteq$) Пусть $a \notin Rad(R)$. Тогда мультипликативное множество $S = \{1, a, a^2, \dots\}$ не содержит нуля . Существует простой идеал $P$, не пересекающийся с $S$. Тогда $a \notin P$, следовательно, $a \notin \bigcap P$ .
	\end{proof}
	
	\newpage
	
	\section{Докажите следующее свойство простых идеалов коммутативного кольца. Если Р - простой идеал в кольце R, и Р содержит пересечение идеалов А, В,..., D кольца R, то один из них содержится в Р. А если Р совпадает с их пересечением, то один из них совпадает с Р.}
	
	\begin{theorem}
		Пусть $P$ — простой идеал кольца $R$. Если $A_1 \cap \dots \cap A_n \subseteq P$, то некоторый идеал $A_i$ содержится в $P$ .
	\end{theorem}
	\begin{proof}
		Предположим противное: ни один $A_i$ не содержится в $P$. Тогда для каждого $i$ существует элемент $a_i \in A_i$, такой что $a_i \notin P$ . Рассмотрим произведение $a_1 a_2 \dots a_n$. Так как идеалы поглощают умножение, это произведение принадлежит пересечению $A_1 \cap \dots \cap A_n$, а значит, принадлежит $P$. Но так как $P$ прост, хотя бы один множитель должен лежать в $P$ — противоречие.
	\end{proof}
	
	\begin{theorem}
		Если $P = A_1 \cap \dots \cap A_n$, то $P = A_i$ для некоторого $i$ .
	\end{theorem}
	\begin{proof}
		Так как $P$ содержит пересечение, по доказанному выше $A_i \subseteq P$ для некоторого $i$. Но так как $P$ равно пересечению всех идеалов, то $P \subseteq A_i$. Следовательно, $P = A_i$.
	\end{proof}
	
	\newpage
	
	\section{Какое кольцо называется факториальным? Какие элементы в факториальном кольце называются простыми? Приведите примеры факториальных колец. Верно ли, что любая область целостности является факториальным кольцом? Ответ обоснуйте примером. Докажите существование НОД и НОК в факториальном кольце.}
	
	\begin{definition}
		Область целостности называется \textbf{факториальной} (ООР), если любой ненулевой необратимый элемент раскладывается в произведение неприводимых, причём это разложение единственно с точностью до порядка и ассоциированности множителей.
	\end{definition}
	
	\begin{definition}
		Ненулевой необратимый элемент $p$ называется \textbf{простым}, если из $p \mid ab$ следует $p \mid a$ или $p \mid b$ .
	\end{definition}
	
	\textbf{Примеры:} Кольца $\mathbb{Z}$ и $F[x]$, где $F$ — поле .
	
	\textbf{Не любая область целостности факториальна.} Контрпример: кольцо $\mathbb{Z}[\sqrt{-5}]$ . В нём элемент 6 имеет два принципиально различных разложения на неприводимые множители: $6 = 2 \cdot 3 = (1+\sqrt{-5})(1-\sqrt{-5})$ .
	
	\begin{theorem}
		В факториальном кольце существуют НОД и НОК любых ненулевых элементов.
	\end{theorem}
	\begin{proof}
		Разложим элементы на простые множители (добавив нулевые степени для выравнивания списка простых делителей): $a = u p_1^{\alpha_1} \dots p_n^{\alpha_n}$, $b = v p_1^{\beta_1} \dots p_n^{\beta_n}$ . \\
		Тогда их НОД конструируется путем выбора минимальных степеней: $(a,b) = p_1^{\min(\alpha_1, \beta_1)} \dots p_n^{\min(\alpha_n, \beta_n)}$. \\
		А НОК конструируется выбором максимальных степеней: $[a,b] = p_1^{\max(\alpha_1, \beta_1)} \dots p_n^{\max(\alpha_n, \beta_n)}$.
	\end{proof}
	
	\newpage
	\setcounter{section}{30}
	
	\section{Докажите, что в факториальном кольце из того, что элемент А взаимно прост с элементами В, С, ..., D следует, что А взаимно прост с их произведением; если элементы В, С, ..., D попарно взаимно просты, и А делится на каждый из них, то А делится на их произведение.}
	
	\begin{theorem}
		Если элемент $A$ взаимно прост с элементами $B, C, \dots, D$, то он взаимно прост с их произведением.
	\end{theorem}
	\begin{proof}
		Пусть некоторый простой делитель $p$ делит одновременно $A$ и произведение $BCD\dots D$. Тогда $p$ делит один из множителей, например $B$. Получаем, что $p$ делит и $A$, и $B$, что ведет к противоречию, так как $(A,B)=1$.
	\end{proof}
	
	\begin{theorem}
		Если элементы $B, C, \dots, D$ попарно взаимно просты и элемент $A$ делится на каждый из них, то $A$ делится на $BCD\dots D$.
	\end{theorem}
	\begin{proof}
		Так как множители попарно взаимно просты, их разложения на простые множители не пересекаются. Поскольку каждый из них делит $A$, все соответствующие простые множители обязательно входят в разложение самого $A$. Следовательно, произведение $BCD\dots D \mid A$.
	\end{proof}
	
	\newpage
	\section{Докажите, что в факториальном кольце R для любого простого элемента р верно: любой элемент а кольца R или делится на р, или взаимно прост с ним; если произведение элементов а, b, c,..., d делится на р, то хотя бы один из них делится на р.}
	
	\begin{theorem}
		Пусть $p$ — простой элемент факториального кольца $R$. Тогда для любого $a \in R$: либо $p \mid a$, либо $(p,a)=1$.
	\end{theorem}
	\begin{proof}
		Если $p$ не делит $a$, то в разложении $a$ на простые множители элемент $p$ не встречается. Следовательно, общих простых делителей у $p$ и $a$ нет, то есть $(p,a)=1$.
	\end{proof}
	
	\begin{theorem}
		Если $p \mid abc\dots d$, то хотя бы один из множителей делится на $p$.
	\end{theorem}
	\begin{proof}
		Так как $p$ прост, из $p \mid ab$ следует $p \mid a$ или $p \mid b$. Применяя это свойство последовательно к произведению нескольких множителей, получаем требуемое утверждение.
	\end{proof}
	
	\newpage
	\section{Докажите, что в кольце главных идеалов любые два элемента имеют НОД. Докажите теорему о представлении НОД в таком кольце.}
	
	\begin{theorem}
		В кольце главных идеалов любые два элемента имеют НОД.
	\end{theorem}
	\begin{proof}
		Пусть $a,b \in R$. Рассмотрим идеал $(a)+(b)$. Так как кольцо является кольцом главных идеалов, $(a)+(b)=(d)$ для некоторого элемента $d$. Так как $a,b \in (d)$, то $d \mid a$ и $d \mid b$, следовательно, $d$ — общий делитель.
		Если $c$ — любой общий делитель $a, b$, то $a,b \in (c)$, значит $(a)+(b) \subseteq (c)$. Отсюда $d \in (c)$, то есть $c \mid d$. Следовательно, $d$ является НОД.
	\end{proof}
	
	\begin{theorem}
		В кольце главных идеалов НОД представим в виде $d = ax+by$.
	\end{theorem}
	\begin{proof}
		Так как $(a)+(b)=(d)$, то $d \in (a)+(b)$. По определению суммы идеалов $d=ax+by$ для некоторых $x,y \in R$.
	\end{proof}
	
	\newpage
	\section{Докажите, что в кольце главных идеалов для любого неприводимого элемента верно: любой элемент кольца или делится на р, или взаимно прост с ним; если произведение нескольких элементов кольца делится на р, то хотя бы один из них делится на р.}
	
	\begin{theorem}
		Пусть $p$ — неприводимый элемент кольца главных идеалов. Тогда для любого $a \in R$: $p \mid a$ или $(p,a)=1$.
	\end{theorem}
	\begin{proof}
		Рассмотрим НОД $d=(p,a)$. Так как $d \mid p$, а $p$ неприводим, то $d \sim p$ или $d \sim 1$. В первом случае $p \mid a$, во втором $(p,a)=1$.
	\end{proof}
	
	\begin{theorem}
		Если $p \mid a_1 a_2 \dots a_n$, то хотя бы один из множителей делится на $p$.
	\end{theorem}
	\begin{proof}
		Достаточно доказать для двух множителей. Пусть $p \mid ab$ и $p \nmid a$. Тогда $(p,a)=1$. Следовательно, существуют $x,y \in R$ такие, что $px+ay=1$. Умножая на $b$, получаем $pb+aby=b$. Так как $p \mid pb$ и $p \mid ab$, то $p \mid b$.
	\end{proof}
	
	\newpage
	\section{Докажите, что в области целостности R любой ненулевой необратимый элемент а, не разлагающийся на неприводимые множители, имеет собственный делитель, который не разлагается на неприводимые множители.}
	
	\begin{theorem}
		Пусть $a$ — ненулевой необратимый элемент области целостности, который не разлагается на неприводимые множители. Тогда $a$ имеет собственный делитель, который также не разлагается на неприводимые множители.
	\end{theorem}
	\begin{proof}
		Так как $a$ не является неприводимым, существуют необратимые элементы $b, c$ такие, что $a=bc$. Если бы оба элемента раскладывались на неприводимые множители, то и $a$ разлагался бы на неприводимые множители, что ведет к противоречию. Следовательно, хотя бы один из элементов $b, c$ не раскладывается на неприводимые множители, и он является собственным делителем элемента $a$.
	\end{proof}
	
	\newpage
	\section{Докажите, что в области целостности R существует бесконечная последовательность элементов, в которой каждый следующий элемент является собственным делителем предыдущего, и существует бесконечная последовательность главных идеалов, в которой каждый следующий строго включает предыдущий.}
	
	\begin{theorem}
		Если в области целостности существует ненулевой необратимый элемент, не разлагающийся на неприводимые множители, то существует бесконечная последовательность $a_1, a_2, a_3, \dots$ такая, что каждый следующий элемент является собственным делителем предыдущего.
	\end{theorem}
	\begin{proof}
		Пусть $a_1 = a$. По предыдущей теореме существует собственный делитель $a_2 \mid a_1$, не разлагающийся на неприводимые множители. Продолжая этот процесс, мы получаем бесконечную последовательность $a_1, a_2, a_3, \dots$, где $a_{n+1} \mid a_n$ и деление является собственным.
	\end{proof}
	
	\begin{theorem}
		Тогда существует бесконечная возрастающая цепочка главных идеалов: $(a_1) \subset (a_2) \subset (a_3) \subset \dots$.
	\end{theorem}
	\begin{proof}
		Так как $a_{n+1} \mid a_n$, то $(a_n) \subseteq (a_{n+1})$. Поскольку делитель является собственным, это включение строгое: $(a_n) \subset (a_{n+1})$.
	\end{proof}
	
	\newpage
	\section{Докажите, что в кольце главных идеалов любая возрастающая последовательность идеалов стабилизируется, а каждый ненулевой необратимый элемент разлагается на неприводимые множители.}
	
	\begin{theorem}
		В кольце главных идеалов любая возрастающая цепочка идеалов стабилизируется.
	\end{theorem}
	\begin{proof}
		Пусть $I_1 \subseteq I_2 \subseteq I_3 \subseteq \dots$ — возрастающая цепочка идеалов. Рассмотрим их объединение $I = \bigcup_{n=1}^\infty I_n$. $I$ является идеалом. Так как кольцо является КГИ, $I = (a)$. Элемент $a \in I_n$ для некоторого $n$. Тогда $I = (a) \subseteq I_n$. Но также $I_n \subseteq I$. Следовательно, $I_n = I$. Значит, для всех $m \ge n$ имеем $I_m = I_n$.
	\end{proof}
	
	\begin{theorem}
		Каждый ненулевой необратимый элемент КГИ раскладывается на неприводимые множители.
	\end{theorem}
	\begin{proof}
		Если бы существовал элемент без такого разложения, то по предыдущим теоремам существовала бы бесконечная строго возрастающая цепочка главных идеалов, что невозможно из-за стабилизации.
	\end{proof}
	
	\newpage
	\section{Докажите факториальность кольца главных идеалов.}
	
	\begin{theorem}
		Любое кольцо главных идеалов является факториальным.
	\end{theorem}
	\begin{proof}
		Существование разложения на неприводимые множители уже доказано. Докажем единственность. Пусть $p$ — неприводимый элемент. Тогда из $p \mid ab$ следует $p \mid a$ или $p \mid b$ (как доказано в 34-м билете). Следовательно, неприводимые элементы в КГИ являются простыми. Теперь стандартным аргументом доказывается единственность разложения на неприводимые множители с точностью до порядка и ассоциированности.
	\end{proof}
	
	\newpage
	\section{Докажите, что евклидово кольцо является кольцом главных идеалов.}
	
	\begin{theorem}
		Любое евклидово кольцо является кольцом главных идеалов.
	\end{theorem}
	\begin{proof}
		Пусть $I$ — ненулевой идеал евклидова кольца. Выберем ненулевой элемент $d \in I$ с минимальным значением евклидовой нормы. Докажем, что $I = (d)$.
		Для любого $a \in I$ разделим его на $d$ с остатком: $a = qd + r$, где $r=0$ или $N(r) < N(d)$. Так как $r = a - qd \in I$, по условию минимальности нормы $d$ случай $N(r) < N(d)$ невозможен. Следовательно, $r=0$. Значит, $a = qd$, откуда $a \in (d)$. Таким образом, $I = (d)$.
	\end{proof}
	
	\newpage
	\section{Докажите, что в кольце главных идеалов каждый простой идеал максимален.}
	
	\begin{theorem}
		В кольце главных идеалов каждый простой идеал максимален.
	\end{theorem}
	\begin{proof}
		Пусть $P$ — простой идеал. Так как кольцо является КГИ, $P = (p)$. Рассмотрим промежуточный идеал $I$, такой что $P \subseteq I \subseteq R$. В КГИ $I = (a)$ для некоторого элемента $a$.
		Из $(p) \subseteq (a)$ следует, что $a \mid p$. Так как $p$ прост и, следовательно, неприводим, имеем $a \sim p$ или $a \sim 1$. 
		В первом случае $I = P$. Во втором случае $I = R$. Следовательно, промежуточных идеалов не существует, и $P$ максимален.
	\end{proof}
	
	\newpage
	\setcounter{section}{40}
	\section{Что такое тело? Приведите пример некоммутативного тела. Докажите, что тело является простым кольцом.}
	
	\begin{definition}
		Тело — это кольцо с единицей, в котором каждый ненулевой элемент обратим.
	\end{definition}
	
	\textbf{Пример некоммутативного тела:} Тело кватернионов $\mathbb{H}$ .
	
	\begin{theorem}
		Любое тело является простым кольцом.
	\end{theorem}
	\begin{proof}
		Пусть $I$ — ненулевой идеал тела $D$. Возьмём $0 \neq a \in I$. Так как $a$ обратим, $a^{-1}a = 1 \in I$. Тогда для любого $x \in D$ имеем $x = x \cdot 1 \in I$. Следовательно, $I = D$. Таким образом, у тела нет нетривиальных идеалов, что по определению означает, что оно простое.
	\end{proof}
	
	\newpage
	\section{Докажите, что ненулевое кольцо с единицей, не содержащее нетривиальных левых идеалов, является телом.}
	
	\begin{theorem}
		Ненулевое кольцо с единицей, не имеющее нетривиальных левых идеалов, является телом.
	\end{theorem}
	\begin{proof}
		Пусть $0 \neq a \in R$. Рассмотрим левый идеал $Ra = \{xa \mid x \in R\}$. Так как $a \neq 0$, идеал ненулевой. По условию в кольце нет нетривиальных идеалов, значит: $Ra = R$. 
		Следовательно, $1 \in Ra$, то есть существует $x \in R$ такой, что $xa = 1$. 
		Аналогично, для элемента $x$ существует $y \in R$ такой, что $ay = 1$. 
		Тогда $x = x(ay) = (xa)y = y$. Следовательно, $xa = ax = 1$. Значит, любой ненулевой элемент обратим, и $R$ является телом.
	\end{proof}
	
	\newpage
	\section{Что такое центр кольца, тела? Каковы его свойства? Приведите примеры.}
	
	\begin{definition}
		Центром кольца (или тела) $R$ называется множество $Z(R) = \{a \in R \mid ax = xa \text{ для всех } x \in R\}$.
	\end{definition}
	
	\begin{theorem}
		Центр кольца является коммутативным подкольцом кольца $R$.
	\end{theorem}
	\begin{proof}
		Если $a, b \in Z(R)$, то для любого $x \in R$ выполнено $(a+b)x = ax+bx = xa+xb = x(a+b)$, следовательно $a+b \in Z(R)$. 
		Аналогично $(ab)x = a(bx) = a(xb) = (ax)b = (xa)b = x(ab)$, поэтому $ab \in Z(R)$. 
		Также $(-a)x = -(ax) = -(xa) = x(-a)$, значит $-a \in Z(R)$. Очевидно, $ab=ba$ для $a,b \in Z(R)$, поэтому подкольцо коммутативно.
	\end{proof}
	
	\textbf{Пример:} Центр кольца матриц $Z(M_n(F)) = \{\lambda E_n \mid \lambda \in F\}$ (скалярные матрицы).
	
	\newpage
	\section{Как определяется структура линейного пространства тела D над подтелом L? Что такое левый базис и левая степень D над L [D:L]? Докажите, что для тел L, D, Е вложенных друг в друга и имеющих конечные левые базисы, [E:L]=[E:D][D:L], и что расширение Е простой степени тела L не имеет собственных подтел, содержащих L.}
	
	\textbf{Структура пространства:} Пусть $L \subseteq D$ — тела. Тело $D$ можно рассматривать как левое линейное пространство над $L$ с операцией умножения на скаляр $\lambda \cdot x = \lambda x$ ($\lambda \in L, x \in D$).
	
	\textbf{Левый базис и степень:} Левым базисом тела $D$ над $L$ называется базис этого линейного пространства $D$. Его мощность (размерность пространства) обозначается $[D:L]$ и называется левой степенью.
	
	\begin{theorem}
		Если $L \subseteq D \subseteq E$ и степени конечны, то $[E:L] = [E:D][D:L]$.
	\end{theorem}
	\begin{proof}
		Пусть $e_1, \dots, e_m$ — базис $E$ над $D$, а $d_1, \dots, d_n$ — базис $D$ над $L$. Элементы $e_i d_j$ образуют систему из $mn$ элементов. Любой $x \in E$ представим как $x = \sum \alpha_i e_i$, где $\alpha_i \in D$. Каждый коэффициент: $\alpha_i = \sum \lambda_{ij} d_j$, $\lambda_{ij} \in L$. Следовательно, $x = \sum \lambda_{ij} (e_i d_j)$. Значит, система $e_i d_j$ порождает $E$. Линейная независимость доказывается аналогично. Поэтому $[E:L] = mn = [E:D][D:L]$.
	\end{proof}
	
	\textbf{Расширение простой степени:} Если степень $[E:L] = p$ (где $p$ — простое число) и $L \subseteq D \subseteq E$, то по доказанной теореме $p = [E:D][D:L]$. Так как $p$ — простое число, один из множителей обязан быть равен 1, а другой $p$. Это означает, что либо $D=L$, либо $D=E$, то есть промежуточных собственных подтел не существует.
	
	\newpage
	\section{Сравнимые элементы по модулю идеала в кольце. Проверить свойства отношения эквивалентности. Каким оно будет в случае нулевого идеала? Описать структуру классов по модулю идеала. Доказать, что если элементы а, b сравнимы, то их степени и кратные тоже сравнимы. Сформулируйте китайскую теорему об остатках для произвольного кольца и для кольца целых чисел.}
	
	\begin{definition}
		Элементы $a, b \in R$ сравнимы по модулю идеала $I$, если $a-b \in I$. Обозначение: $a \equiv b \pmod I$.
	\end{definition}
	В случае нулевого идеала $I=(0)$ сравнимость совпадает с обычным равенством.
	
	\begin{theorem}
		Отношение сравнимости по модулю идеала является отношением эквивалентности.
	\end{theorem}
	\begin{proof}
		Рефлексивность: $a-a=0 \in I$. Симметричность: если $a-b \in I$, то $b-a = -(a-b) \in I$. Транзитивность: если $a-b \in I$ и $b-c \in I$, то $a-c = (a-b)+(b-c) \in I$.
	\end{proof}
	
	\textbf{Структура классов:} Классом вычетов элемента $a$ является смежный класс $[a] = a+I = \{a+x \mid x \in I\}$.
	
	\begin{theorem}
		Если $a \equiv b \pmod I$, то $a^n \equiv b^n \pmod I$ и $ra \equiv rb \pmod I$.
	\end{theorem}
	\begin{proof}
		Для степеней: $a^n - b^n = (a-b)(a^{n-1} + \dots + b^{n-1}) \in I$. Для кратных: $ra - rb = r(a-b) \in I$.
	\end{proof}
	
	\begin{theorem}[Китайская теорема об остатках]
		Пусть идеалы $I_1, \dots, I_n$ попарно взаимно просты. Тогда отображение $R \to R/I_1 \times \dots \times R/I_n$ является сюръективным, а его ядро равно $I_1 \cap \dots \cap I_n$. (Для $\mathbb{Z}$: если числа $m_1 \dots m_n$ попарно взаимно просты, то система сравнений $x \equiv c_k \pmod{m_k}$ всегда имеет решение).
	\end{theorem}
	
	\newpage
	\section{Доказать, что в коммутативном кольце R с единицей идеал Р простой тогда и только тогда, когда R/P --- область целостности. Идеал М --- максимальный тогда и только тогда, когда R / M --- поле. Сформулируйте последнее утверждение для кольца многочленов.}
	
	\begin{theorem}
		Идеал $P$ прост $\iff$ $R/P$ — область целостности.
	\end{theorem}
	\begin{proof}
		В факторкольце $(a+P)(b+P) = P \iff ab \in P$. Так как $P$ прост, $a \in P$ или $b \in P$. Следовательно, $a+P=P$ или $b+P=P$. Делителей нуля нет.
	\end{proof}
	
	\begin{theorem}
		Идеал $M$ максимален $\iff R/M$ — поле.
	\end{theorem}
	\begin{proof}
		Пусть $M$ максимален и $a+M \neq M$. Тогда $a \notin M$. Из максимальности $M+(a)=R$. Значит, существуют $m \in M, r \in R$ такие, что $m+ra=e$. Тогда $(r+M)(a+M) = e+M$, то есть каждый ненулевой элемент обратим. Обратно, если $R/M$ — поле, факторкольцо не имеет нетривиальных идеалов. Следовательно, между $M$ и $R$ нет промежуточных идеалов, значит $M$ максимален.
	\end{proof}
	
	\textbf{Для многочленов:} В кольце $F[x]$ над полем $F$ идеал $(f)$ максимален (и, соответственно, факторкольцо является полем) тогда и только тогда, когда многочлен $f$ неприводим.
	
	\newpage
	\section{Дайте определение гомоморфизма колец. Доказать, что гомоморфный образ кольца является кольцом. Является ли гомоморфный образ области целостности областью целостности? Верно ли, что гомоморфный образ кольца с делителями нуля может быть областью целостности?}
	
	\begin{definition}
		Отображение $f: R \to S$ называется гомоморфизмом колец, если $f(a+b)=f(a)+f(b)$ и $f(ab)=f(a)f(b)$.
	\end{definition}
	
	\begin{theorem}
		Гомоморфный образ кольца является кольцом.
	\end{theorem}
	\begin{proof}
		Множество образов $f(R) = \{f(a) \mid a \in R\}$ замкнуто относительно сложения и умножения: $f(a)+f(b)=f(a+b)$, $f(a)f(b)=f(ab)$, а также содержит противоположные элементы $-f(a)=f(-a)$. Следовательно, $f(R)$ является кольцом.
	\end{proof}
	
	\textbf{Образ области целостности:} Гомоморфный образ области целостности не обязан быть областью целостности. Пример: гомоморфизм $\mathbb{Z} \to \mathbb{Z}/6\mathbb{Z}$.
	
	\textbf{Образ кольца с делителями нуля:} Да, верно. Гомоморфный образ кольца с делителями нуля может быть областью целостности. Пример: гомоморфизм $\mathbb{Z}/6\mathbb{Z} \to \mathbb{Z}/3\mathbb{Z}$.
	
	\newpage
	\section{Докажите, что ядром кольцевого гомоморфизма является идеал. Докажите основную теорему о гомоморфизме колец. Что можно сказать о гомоморфизме с нулевым ядром?}
	
	\begin{definition}
		Ядром гомоморфизма $f: R \to S$ называется множество $\ker f = \{a \in R \mid f(a)=0\}$.
	\end{definition}
	
	\begin{theorem}
		Ядро гомоморфизма является идеалом.
	\end{theorem}
	\begin{proof}
		Если $a, b \in \ker f$, то $f(a-b)=f(a)-f(b)=0$, следовательно $a-b \in \ker f$. Для любого $r \in R$, $f(ra)=f(r)f(a)=0$, поэтому $ra \in \ker f$.
	\end{proof}
	
	\begin{theorem}[Основная теорема о гомоморфизме]
		Для любого гомоморфизма колец $f: R \to S$ имеем изоморфизм $R/\ker f \cong f(R)$.
	\end{theorem}
	\begin{proof}
		Определим отображение $\varphi: R/\ker f \to f(R)$ по правилу $\varphi(a+\ker f) = f(a)$. Оно корректно, так как из $a-b \in \ker f$ следует $f(a)=f(b)$. Отображение является сюръективным по определению образа. Ядро этого отображения тривиально, следовательно, оно инъективно и является изоморфизмом.
	\end{proof}
	
	\textbf{Нулевое ядро:} Если $\ker f = (0)$, то гомоморфизм инъективен.
	
	\newpage
	\section{Пусть f гомоморфизм коммутативного кольца R с единицей в кольцо S. Докажите, что если Р --- простой идеал в S, то его прообраз --- простой идеал в R. Прообраз максимального идеала прост, а если f --- сюръекция, то прообраз максимального идеала максимален.}
	
	\begin{theorem}
		Если $P$ — простой идеал в $S$, то его прообраз $f^{-1}(P)$ является простым идеалом в $R$.
	\end{theorem}
	\begin{proof}
		Пусть $ab \in f^{-1}(P)$. Тогда $f(ab) = f(a)f(b) \in P$. Так как $P$ прост, $f(a) \in P$ или $f(b) \in P$. Следовательно, $a \in f^{-1}(P)$ или $b \in f^{-1}(P)$.
	\end{proof}
	
	\begin{theorem}
		Прообраз максимального идеала при сюръективном гомоморфизме является максимальным.
	\end{theorem}
	\begin{proof}
		Пусть $M \subseteq S$ максимальный идеал и $f$ сюръективен. По теоремам об изоморфизме имеем $R/f^{-1}(M) \cong S/M$. Так как $S/M$ — поле (поскольку $M$ максимален), то и $R/f^{-1}(M)$ является полем. Следовательно, идеал $f^{-1}(M)$ максимален.
	\end{proof}
	
	\newpage
	\section{Как определяется характеристика области целостности? Как выглядит формула бинома Ньютона в области целостности простой характеристики? Какое поле называется простым? Доказать, что поле вычетов и поле рациональных чисел просты.}
	
	\begin{definition}
		Характеристикой области целостности называется минимальное натуральное число $n$ такое, что $n \cdot 1 = 0$. Если такого числа нет, характеристика равна нулю. (Характеристика всегда равна либо 0, либо простому числу).
	\end{definition}
	
	\begin{theorem}
		В области целостности характеристики $p$ выполняется формула: $(a+b)^p = a^p + b^p$.
	\end{theorem}
	\begin{proof}
		По классической формуле бинома Ньютона: $(a+b)^p = \sum_{k=0}^p \binom{p}{k} a^k b^{p-k}$. Для всех $1 \le k \le p-1$ биномиальный коэффициент $\binom{p}{k}$ делится на $p$, поэтому в кольце характеристики $p$ соответствующие слагаемые равны нулю.
	\end{proof}
	
	\begin{definition}
		Поле называется \textbf{простым}, если оно не содержит собственных подполей.
	\end{definition}
	
	\begin{theorem}
		Поле рациональных чисел $\mathbb{Q}$ и поле вычетов $\mathbb{F}_p$ являются простыми.
	\end{theorem}
	\begin{proof}
		Любое подполе по определению должно содержать единицу $1$, а значит, в силу замкнутости по сложению, и все элементы вида $n \cdot 1$. Следовательно, в поле характеристики 0 подполе обязано содержать все рациональные дроби $\mathbb{Q}$, а в поле характеристики $p$ оно содержит всё поле $\mathbb{F}_p$.
	\end{proof}
	
	\newpage
	\section{Как определяются внешняя и внутренняя прямые суммы колец? Проекторы и идеалы в случае внешней прямой суммы. Как устроена группа обратимых элементов внешней прямой суммы колец?}
	
	\textbf{Внешняя прямая сумма:} Внешней прямой суммой колец $R_1 \oplus R_2$ называется их декартово произведение $R_1 \times R_2$, снабженное покомпонентно заданными операциями сложения и умножения: $(a, b) + (c, d) = (a+c, b+d)$ и $(a, b) \cdot (c, d) = (ac, bd)$.
	
	\textbf{Внутренняя прямая сумма:} Говорят, что кольцо $R$ разлагается во внутреннюю прямую сумму $R = I_1 \oplus I_2$ своих двусторонних идеалов $I_1$ и $I_2$, если $I_1 \cap I_2 = \{0\}$ и $I_1 + I_2 = R$.
	
	\textbf{Проекторы и идеалы:} Во внешней прямой сумме элементами $e_1 = (1_R, 0)$ и $e_2 = (0, 1_R)$ являются идемпотенты (проекторы), причем они ортогональны: $e_1 e_2 = 0$. Любой идеал во внешней прямой сумме имеет вид $J_1 \oplus J_2$, где $J_1$ — идеал в $R_1$, а $J_2$ — идеал в $R_2$.
	
	\textbf{Группа обратимых элементов:} Элемент $(u_1, u_2)$ обратим в $R_1 \oplus R_2$ тогда и только тогда, когда существует элемент $(v_1, v_2)$ такой, что $(u_1 v_1, u_2 v_2) = (1, 1)$. Это эквивалентно тому, что $u_1$ обратим в $R_1$ и $u_2$ обратим в $R_2$. Таким образом, мультипликативная группа обратимых элементов $U(R_1 \oplus R_2)$ изоморфна прямому произведению групп $U(R_1) \times U(R_2)$.
	
	\newpage
	\section{Многочлены с одной переменной, полиномиальные функции и подстановка в многочлен. Возможна ли ситуация, когда подстановка в ненулевой многочлен приводит к нулевой функции? Докажите, что поле комплексных чисел изоморфно фактор кольцу кольца многочленов над полем действительных чисел.}
	
	\textbf{Полиномиальные функции:} Для многочлена $f(x) \in R[x]$ полиномиальной функцией называется отображение $a \mapsto f(a)$, получаемое подстановкой элемента $a \in R$ вместо формальной переменной $x$.
	
	\textbf{Ненулевой многочлен и нулевая функция:} Да, такая ситуация возможна в конечных кольцах. Например, над конечным полем $\mathbb{F}_p$ многочлен $x^p - x$ не является нулевым как формальный многочлен (так как его коэффициенты не равны нулю), но он обращается в нуль при подстановке абсолютно любого элемента из поля $\mathbb{F}_p$ (в силу малой теоремы Ферма), то есть задает нулевую функцию.
	
	\begin{theorem}
		Поле комплексных чисел изоморфно факторкольцу: $\mathbb{C} \cong \mathbb{R}[x]/(x^2+1)$.
	\end{theorem}
	\begin{proof}
		Рассмотрим гомоморфизм $\varphi: \mathbb{R}[x] \to \mathbb{C}$, заданный правилом $\varphi(f) = f(i)$ (подстановка мнимой единицы). Он сюръективен, поскольку любое комплексное число $a+bi$ можно получить как образ линейного многочлена: $a+bi = \varphi(a+bx)$. 
		Ядро этого гомоморфизма состоит из многочленов, которые обращаются в нуль в точке $i$, то есть делятся на $x^2+1$ (так как $i^2+1=0$). Следовательно, $\ker \varphi = (x^2+1)$. 
		По основной теореме о гомоморфизме получаем требуемый изоморфизм: $\mathbb{R}[x]/(x^2+1) \cong \mathbb{C}$.
	\end{proof}
	
	\newpage
	\setcounter{section}{52}
	\section{Какой элемент называется алгебраическим над полем? Докажите эквивалентность трёх свойств алгебраического элемента.}
	
	\begin{definition}
		Элемент $\alpha$ расширения поля $E$ над полем $K$ называется \textbf{алгебраическим} над $K$, если существует ненулевой многочлен $f(x) \in K[x]$ такой, что $f(\alpha)=0$.
	\end{definition}
	
	\begin{theorem}
		Для элемента $\alpha$ эквивалентны следующие 3 условия:
		1. $\alpha$ алгебраичен над $K$.
		2. Множество $1, \alpha, \alpha^2, \dots$ линейно зависимо над $K$.
		3. Поле $K(\alpha)$ имеет конечную степень над $K$.
	\end{theorem}
	\begin{proof}
		$(1 \Rightarrow 2)$: Если $a_0 + a_1\alpha + \dots + a_n\alpha^n = 0$ (где не все коэффициенты нули, так как $\alpha$ корень ненулевого многочлена), то степени линейно зависимы по определению. \\
		$(2 \Rightarrow 3)$: Если степени зависимы, то все высокие степени $\alpha$ линейно выражаются через конечное число меньших степеней. Поэтому пространство $K(\alpha)$ порождается конечным числом элементов, то есть имеет конечную размерность (степень) над $K$. \\
		$(3 \Rightarrow 1)$: Конечномерное пространство не может содержать бесконечно много линейно независимых элементов. Значит, степени $1, \alpha, \alpha^2, \dots$ линейно зависимы, что дает многочлен, зануляющийся в $\alpha$.
	\end{proof}
	
	\newpage
	\section{Какой элемент называется трансцендентным над полем? Докажите эквивалентность трёх свойств трансцендентного элемента.}
	
	\begin{definition}
		Элемент $\alpha$ расширения поля $E$ над полем $K$ называется \textbf{трансцендентным} над $K$, если не существует ненулевого многочлена $f(x) \in K[x]$ такого, что $f(\alpha)=0$.
	\end{definition}
	
	\begin{theorem}
		Для элемента $\alpha$ эквивалентны следующие 3 свойства:
		1. $\alpha$ трансцендентен над $K$.
		2. Система $1, \alpha, \alpha^2, \dots$ линейно независима над $K$.
		3. Поле $K(\alpha)$ изоморфно полю рациональных дробей $K(x)$.
	\end{theorem}
	\begin{proof}
		$(1 \Rightarrow 2)$: Если бы степени были линейно зависимы, то существовал бы ненулевой многочлен, обращающийся в нуль на $\alpha$, что противоречит трансцендентности. \\
		$(2 \Rightarrow 3)$: Отображение подстановки $K[x] \to K[\alpha]$, $f(x) \mapsto f(\alpha)$ инъективно, так как ядро тривиально (система линейно независима). Поэтому $K[\alpha] \cong K[x]$. Переходя к полям частных, получаем $K(\alpha) \cong K(x)$. \\
		$(3 \Rightarrow 1)$: В поле $K(x)$ формальная переменная $x$ не удовлетворяет никакому ненулевому многочлену над $K$. В силу изоморфизма, то же верно и для $\alpha$.
	\end{proof}
	
	\newpage
	\section{Минимальный многочлен алгебраического элемента. Его свойство. Степень алгебраического элемента над полем и подполем. Алгебраические элементы первой степени.}
	
	\begin{definition}
		\textbf{Минимальным многочленом} алгебраического элемента $\alpha$ над полем $K$ называется неприводимый унитарный многочлен $m_\alpha(x) \in K[x]$ минимальной степени, такой что $m_\alpha(\alpha)=0$.
	\end{definition}
	
	\begin{theorem}[Свойство минимального многочлена]
		Если $f(\alpha)=0$, то минимальный многочлен делит $f$: $m_\alpha(x) \mid f(x)$.
	\end{theorem}
	\begin{proof}
		Разделим $f$ на $m_\alpha$ с остатком: $f = q \cdot m_\alpha + r$, где $\deg r < \deg m_\alpha$. Подставляя $\alpha$, получаем: $0 = f(\alpha) = q(\alpha)m_\alpha(\alpha) + r(\alpha) = r(\alpha)$. В силу минимальности степени $m_\alpha$, остаток $r$ обязан быть нулевым: $r=0$. Следовательно, $m_\alpha \mid f$.
	\end{proof}
	
	\textbf{Степень:} Степенью алгебраического элемента $\alpha$ называется размерность $[K(\alpha):K]$. Она в точности равна степени его минимального многочлена, так как базис $K(\alpha)$ над $K$ образуют элементы $1, \alpha, \dots, \alpha^{n-1}$, где $n = \deg m_\alpha$.
	
	\textbf{Первая степень:} Алгебраические элементы первой степени — это в точности сами элементы базового поля: $\alpha \in K$.
	
	\newpage
	\section{Неприводимые многочлены в R[x] и С[х]. Докажите, что над произвольным полем существует бесконечно много унитарных неприводимых многочленов. Докажите, что над конечным полем существуют унитарные неприводимые многочлены сколь угодно большой степени.}
	
	\begin{theorem}
		Неприводимые многочлены над $\mathbb{C}$ имеют строго степень 1. Неприводимые многочлены над $\mathbb{R}$ имеют степень 1 или 2 (причем квадратичные имеют отрицательный дискриминант).
	\end{theorem}
	\begin{proof}
		По основной теореме алгебры любой многочлен над $\mathbb{C}$ имеет корень $\lambda$, а значит отщепляет множитель $(x-\lambda)$, поэтому неприводимы только линейные. Комплексные корни вещественного многочлена возникают сопряженными парами $a \pm bi$. Их произведение дает неприводимый над $\mathbb{R}$ квадратичный множитель $x^2+px+q$.
	\end{proof}
	
	\begin{theorem}
		Над любым полем существует бесконечно много унитарных неприводимых многочленов.
	\end{theorem}
	\begin{proof}
		Предположим противное: пусть $f_1, \dots, f_n$ — вообще все неприводимые унитарные многочлены. Рассмотрим многочлен $F = f_1 f_2 \dots f_n + 1$. Многочлен $F$ при делении на любой из $f_i$ дает остаток 1, то есть не делится ни на один из них. Противоречие.
	\end{proof}
	
	\begin{theorem}
		Над конечным полем существуют неприводимые многочлены сколь угодно большой степени.
	\end{theorem}
	\begin{proof}
		Если бы степени были ограничены некоторым числом, то (поскольку поле конечно) существовало бы лишь конечное число неприводимых многочленов. Но тогда, перемножая их, мы получили бы лишь конечное число всех возможных многочленов в кольце, что очевидно неверно. (Альтернативно: всегда существует расширение $\mathbb{F}_{q^d}$, порождаемое примитивным элементом $\alpha$, чей минимальный многочлен имеет в точности степень $d$).
	\end{proof}
	
	\newpage
	\section{Примитивные многочлены над факториальным кольцом. Теорема о представлении произвольного многочлена в виде произведения примитивного и рациональной дроби. Лемма Гаусса с доказательством.}
	
	\begin{definition}
		Многочлен над факториальным кольцом называется \textbf{примитивным}, если наибольший общий делитель (НОД) всех его коэффициентов равен 1.
	\end{definition}
	
	\begin{theorem}
		Любой многочлен $f(x) \in K[x]$ можно представить в виде $f(x) = c \cdot g(x)$, где $c \in K$ (рациональная дробь/константа), а $g(x)$ — примитивный многочлен.
	\end{theorem}
	
	\begin{theorem}[Лемма Гаусса]
		Произведение примитивных многочленов является примитивным многочленом.
	\end{theorem}
	\begin{proof}
		Пусть $f, g$ — примитивные многочлены. Предположим, что их произведение $fg$ непримитивно. Тогда существует простой элемент $p$, который делит абсолютно все коэффициенты произведения $fg$. \\
		Так как $f$ и $g$ примитивны, $p$ не делит все их коэффициенты. Выберем в многочленах $f$ и $g$ первые с конца (старшие из неделящихся) коэффициенты, которые \textit{не делятся} на $p$. В коэффициенте произведения $fg$ при соответствующей суммарной степени появится слагаемое, равное произведению этих двух неделящихся коэффициентов, плюс другие слагаемые, которые уже гарантированно кратны $p$. В итоге этот конкретный коэффициент произведения $fg$ не будет делиться на $p$. Полученное противоречие доказывает лемму.
	\end{proof}
	
	\newpage
	\section{Неприводимость над $\mathbb{Z}$ и над $\mathbb{Q}$. Критерий Эйзенштейна. Редукционный критерий.}
	
	\begin{theorem}
		Многочлен из $\mathbb{Z}[x]$ примитивен и неприводим над кольцом целых чисел $\mathbb{Z}$ тогда и только тогда, когда он неприводим над полем рациональных чисел $\mathbb{Q}$.
	\end{theorem}
	\begin{proof}
		Если многочлен раскладывается над $\mathbb{Q}$, мы можем вынести общие знаменатели дробей. После приведения к общему знаменателю и применения леммы Гаусса (которая позволяет сократить лишние константы), мы получим разложение этого же многочлена на множители уже над $\mathbb{Z}$.
	\end{proof}
	
	\begin{theorem}[Критерий Эйзенштейна]
		Пусть дан многочлен $f(x) = a_n x^n + \dots + a_0 \in \mathbb{Z}[x]$. Если существует простое число $p$ такое, что:
		1. $p \mid a_0, \dots, p \mid a_{n-1}$ (все коэффициенты кроме старшего делятся на $p$),
		2. $p \nmid a_n$ (старший коэффициент не делится на $p$),
		3. $p^2 \nmid a_0$ (свободный член не делится на $p^2$),
		то многочлен $f(x)$ неприводим над $\mathbb{Q}$.
	\end{theorem}
	
	\begin{theorem}[Редукционный критерий]
		Если многочлен из $\mathbb{Z}[x]$ после взятия его коэффициентов по модулю простого числа $p$ (редукции) остается неприводимым над полем вычетов $\mathbb{F}_p$ (и его степень при этом не падает), то исходный многочлен неприводим над $\mathbb{Q}$.
	\end{theorem}
	
	\newpage
	\section{Если кольцо R факториально, то R[x] --- тоже.}
	
	\begin{theorem}
		Если кольцо $R$ факториально (является областью с однозначным разложением), то кольцо многочленов $R[x]$ также факториально.
	\end{theorem}
	\begin{proof}
		Доказательство опирается на поле частных кольца $R$ и Лемму Гаусса. Каждый ненулевой многочлен $f \in R[x]$ можно вынести за скобки НОД своих коэффициентов (содержимое), представив в виде $f = c \cdot f_0$, где $f_0$ — примитивный многочлен. \\
		В поле частных $Frac(R)[x]$ многочлен $f_0$ однозначно раскладывается на неприводимые множители (так как над полем кольцо многочленов является КГИ и факториально). С помощью леммы Гаусса мы можем избавиться от знаменателей и перенести это разложение обратно в кольцо $R[x]$. Уникальность разложения коэффициента $c$ следует из факториальности исходного кольца $R$. Таким образом, существование и единственность разложения сохраняются.
	\end{proof}
	
	\newpage
	\section{Многочлены от нескольких переменных. Степень, однородные многочлены. Если R --- область целостности, то и R[x,y,...,z] --- тоже. Факториальность. Существование неприводимых многочленов любой степени над любым полем.}
	
	\textbf{Определения:} Кольцо многочленов от $n$ переменных определяется индуктивно: $R[x_1, \dots, x_n] := (R[x_1, \dots, x_{n-1}])[x_n]$. \textbf{Степенью одночлена} называется сумма его показателей $\sum i_k$. Степенью многочлена — максимум степеней его одночленов. Многочлен \textbf{однороден}, если все его ненулевые одночлены имеют одинаковую степень.
	
	\begin{theorem}
		Если $R$ — область целостности, то $R[x_1, \dots, x_n]$ также область целостности.
	\end{theorem}
	\begin{proof}
		Индукция по $n$. База $n=1$: старший коэффициент произведения ненулевых многочленов равен произведению их старших коэффициентов. Так как делителей нуля нет в $R$, произведение не ноль. Переход: $R[x_1, \dots, x_n] \cong S[x_n]$, где $S = R[x_1, \dots, x_{n-1}]$. По предположению индукции $S$ — область целостности, следовательно $S[x_n]$ тоже.
	\end{proof}
	
	\textbf{Факториальность:} Если $R$ факториально, то и $R[x_1, \dots, x_n]$ факториально. Доказывается аналогичной индукцией по теореме Гаусса ($A$ факториально $\implies A[x]$ факториально).
	
	\textbf{Неприводимые многочлены:} Над любым полем $K$ существуют неприводимые многочлены любой степени $d \ge 1$. Для $\mathbb{Q}$ — это многочлены $x^d-2$ по Эйзенштейну. Для конечного поля $\mathbb{F}_q$ — это минимальный многочлен элемента $\alpha$, порождающего единственное расширение поля $\mathbb{F}_{q^d}$, его степень в точности $d$.
	
	\newpage
	\section{Доказать, что R[x,y,...,z], где R --- поле, не является кольцом главных идеалов.}
	
	\begin{theorem}
		Кольцо многочленов от двух и более переменных $K[x, y]$ над полем $K$ не является кольцом главных идеалов.
	\end{theorem}
	\begin{proof}
		Рассмотрим идеал $I = \langle x, y \rangle$, порожденный переменными $x$ и $y$ (множество многочленов без свободного члена). Предположим, что $I$ главный, то есть существует многочлен $d(x,y)$, такой что $I = \langle d(x,y) \rangle$. \\
		Тогда $x \in \langle d \rangle \implies d \mid x$. Отсюда степень $d$ по $y$ равна 0, а по $x$ не больше 1. \\
		Аналогично, $y \in \langle d \rangle \implies d \mid y$, откуда степень $d$ по $x$ равна 0. \\
		Объединяя, получаем, что $d(x,y) = c \in K^\times$ (константа). Но любой идеал, порожденный константой, совпадает со всем кольцом: $\langle c \rangle = K[x,y]$. В то же время $1 \notin I$, значит $I \neq K[x,y]$. Противоречие. Идеал $\langle x,y \rangle$ не главный.
	\end{proof}
	
	\newpage
	\section{Симметрические многочлены. Алгебраически зависимые элементы над кольцом R. Независимость элементарных симметрических многочленов. Основная теорема о симметрических многочленах.}
	
	\begin{definition}
		Многочлен $f(x_1, \dots, x_n)$ называется \textbf{симметрическим}, если он не меняется при любой перестановке своих переменных.
	\end{definition}
	
	\textbf{Элементарные симметрические многочлены:}
	Задаются как коэффициенты многочлена $(t-x_1)\dots(t-x_n)$ (формулы Виета):
	$\sigma_1 = \sum x_i$, $\sigma_2 = \sum_{i<j} x_i x_j$, $\dots$, $\sigma_n = x_1 x_2 \dots x_n$. Они алгебраически независимы над кольцом, что означает, что не существует ненулевого многочлена $P$, зануляющегося при подстановке $\sigma_1, \dots, \sigma_n$.
	
	\begin{theorem}[Основная теорема о симметрических многочленах]
		Любой симметрический многочлен единственным образом выражается в виде многочлена от элементарных симметрических многочленов $\sigma_1, \dots, \sigma_n$.
	\end{theorem}
	\begin{proof}
		Используется лексикографический порядок старших членов. Берется старший член исходного симметрического многочлена, и для него подбирается такая комбинация $\sigma_1^{k_1}\dots\sigma_n^{k_n}$, старший член которой в точности с ним совпадает. Затем эта комбинация вычитается из исходного многочлена. Старший член остатка в лексикографическом порядке строго меньше. Процесс повторяется и обрывается за конечное число шагов.
	\end{proof}
	
	\newpage
	\section{Теорема Кронеккера с доказательством.}
	
	\begin{theorem}[Теорема Кронеккера]
		Для любого многочлена $f(x) \in K[x]$ существует расширение поля, в котором этот многочлен раскладывается на линейные множители.
	\end{theorem}
	\begin{proof}
		Пусть многочлен $f(x)$ имеет степень $n$. Доказательство ведется индукцией по $n$.
		Если многочлен $f$ имеет неприводимый нелинейный делитель $p(x) \in K[x]$, мы можем искусственно присоединить к полю $K$ его корень, построив факторкольцо $E_1 = K[t]/(p(t))$. В этом новом поле класс вычетов $[t]$ будет корнем многочлена $p(x)$, а значит и $f(x)$.
		Таким образом, в поле $E_1$ многочлен отщепляет линейный множитель: $f(x) = (x-\alpha)g(x)$. 
		Далее мы применяем тот же процесс для многочлена $g(x)$. Так как его степень меньше, через конечное число шагов мы построим искомое поле разложения, где $f(x)$ распадется на линейные множители.
	\end{proof}
	
	\newpage
	\section{Задача об общих корнях двух многочленов над полем К. Результант.}
	
	\begin{definition}
		\textbf{Результантом} двух многочленов $f$ и $g$ (обозначается $Res(f,g)$) называется определитель их матрицы Сильвестра (составленной из коэффициентов $f$ и $g$ по специальному правилу смещения).
	\end{definition}
	
	\begin{theorem}
		Многочлены $f$ и $g$ имеют общий корень (в некотором расширении поля) тогда и только тогда, когда $Res(f,g) = 0$.
	\end{theorem}
	\begin{proof}
		Если $f$ и $g$ имеют общий корень $\alpha$, то $f(\alpha)=g(\alpha)=0$. Если подставить $\alpha$ в матричное уравнение, связывающее матрицу Сильвестра со степенями $\alpha$, мы получим нулевой вектор. Существование нетривиального решения однородной системы означает, что определитель матрицы равен нулю. \\
		Обратно, если результант равен нулю (определитель матрицы Сильвестра нулевой), то строки матрицы линейно зависимы. Это алгебраически означает существование ненулевых многочленов $u, v$ таких, что $\deg u < \deg g$, $\deg v < \deg f$ и $u f + v g = 0$. Отсюда следует, что $f$ и $g$ не взаимно просты (иначе $u f + v g$ порождало бы единицу). А так как у них есть нетривиальный общий делитель, у них есть и общий корень в поле разложения этого делителя.
	\end{proof}

\end{document}